\section{Normalizability and gravity-decoupling limits}
\label{sec:decoupling}

The compact theorem and a rigid noncompact limit impose different
normalizability conditions. In the compact model, the rows of $Q_X$ arise from
dynamical vector fields. In a noncompact limit, some of those fields can
become background flavor sources, in which case their moment-map equations
must be removed. We formulate the distinction in this section.

\subsection{Finite compact volume}
\label{subsec:compact-normalizability}

On a smooth compact resolution, the standard reduction
\eqref{eq:c3-reduction} gives the five-dimensional kinetic term
\begin{equation}
 S_5\supset-\frac14\int_{M_5}G_{IJ}F^I\wedge *_5F^J,
 \qquad
 G_{IJ}\propto\int_{\wideX}\omega_I\wedge *_6\omega_J.
 \label{eq:vector-kinetic-matrix}
\end{equation}
\cite{Cadavid:1995bk}.
Every harmonic form on a compact manifold has finite norm. After separating
the graviphoton, these fields supply the $h^{1,1}(\wideX)-1$ vector
multiplets. At the transition point the graviphoton direction is the
K\"ahler class of the singular threefold, which pairs to zero with every
exceptional curve and root, so its row of $Q_X$ vanishes and the effective
rows in Section~\ref{sec:examples} are vector-multiplet directions. The
pairings in \eqref{eq:nodal-charge} and
\eqref{eq:cocharacter} are therefore charges under dynamical fields, not
charges under nondynamical flavor backgrounds.

Similarly, the kinetic metric for compact complex-structure moduli is
schematically
\begin{equation}
 G_{\kappa\bar\lambda}
 \propto
 \frac{\int_X\chi_\kappa\wedge\overline\chi_{\bar\lambda}}
      {\int_X\Omega\wedge\overline\Omega},
 \qquad
 \chi_\kappa\in H^{2,1}(X).
 \label{eq:complex-structure-metric}
\end{equation}
\cite{Cadavid:1995bk}.
The local coordinates in $\cB_\sigma$ are limits of these compact directions,
subject to the global image theorem. Thus no additional local $L^2$
assumption is needed to impose
\begin{equation}
 Q_Xb=0
 \label{eq:full-compact-constraint}
\end{equation}
at finite compact volume.

\subsection{Dynamical and background vectors in a local limit}
\label{subsec:dynamic-background}

Consider a one-parameter family of K\"ahler classes $J(\Lambda)$ with
\begin{equation}
 \Vol_{J(\Lambda)}(X)\longrightarrow\infty,
 \qquad
 G_{N,5}\longrightarrow0.
 \label{eq:gravity-decoupling}
\end{equation}
For a vector combination $v=v^ID_I$, let
$\omega_v=v^I\omega_I$ and define its kinetic norm
\begin{equation}
 \mathcal N_v(\Lambda)
 =\int_X\omega_v\wedge*_{J(\Lambda)}\omega_v,
 \label{eq:vector-norm}
\end{equation}
understood in a finite-normalization basis after the standard five-dimensional
Einstein-frame rescaling. The vector belongs to the retained dynamical theory
only when its normalized kinetic term has a finite nonzero limit. Geometrically,
this is the requirement that it become an $L^2$ two-form on the retained
noncompact space. A compact divisor may furnish such a form; a divisor which
becomes noncompact gives a non-normalizable flavor background.

Let $P_{\mathrm{dyn}}$ project the compact vector lattice onto the combinations
that remain dynamical, and let $\cB_{\mathrm{loc}}$ be the product of the
local invariant-coordinate spaces retained in the chosen noncompact limit.
The moment-map matrix of the limiting theory is
\begin{equation}
 \boxed{
 Q_{\mathrm{dyn}}=P_{\mathrm{dyn}}Q_X,
 \qquad
 \operatorname{Im}(\cM_H\to\cB_{\mathrm{loc}})
 =\ker Q_{\mathrm{dyn}}.}
 \label{eq:projected-decoupling-matrix}
\end{equation}
If the norm of a shared vector diverges, its canonically normalized coupling
to the local sector vanishes. That vector is no longer integrated over in the
local theory; its moment-map equation is not imposed, and the corresponding
row of $Q_X$ is deleted. Equation \eqref{eq:projected-decoupling-matrix} is
therefore the limit of the compact derivation, not an independent ansatz.

\subsection{Multistage decoupling}
\label{subsec:multistage}

Several local singularities can remain connected after the overall
Calabi--Yau volume is sent to infinity. Shared compact surfaces or
lower-codimension singular loci can continue to gauge a common flavor subgroup.
Only a further decompactification factorizes the individual local theories.
This multistage structure appears in geometric constructions of 5d theories
and in compact analyses of generalized symmetries
\cite{Apruzzi:2019vpe,Cvetic:2023plv}.

For the present mechanism there are accordingly three regimes:
\begin{equation}
\begin{array}{c|c|c}
\text{regime}&\text{dynamical shared vectors}&\text{constraint}\\
\hline
\text{compact supergravity}&\text{all compact rows}&Q_Xb=0\\
\text{glued local QFT}&\text{selected finite-norm rows}&P_{\mathrm{dyn}}Q_Xb=0\\
\text{factorized local sectors}&\text{no shared rows}&\text{no cross-sector equation}.
\end{array}
\label{eq:three-regimes}
\end{equation}
Coupling to gravity is therefore sufficient but not necessary for the
correlation. A gravity-decoupled glued QFT may retain it as an ordinary 5d
gauging. What is special to the compact statement is that every row is
dynamical before a limit is chosen.

\subsection[Complex deformations and the local L2 caveat]{Complex deformations and the local $L^2$ caveat}
\label{subsec:local-l2-caveat}

The same distinction applies to Higgs coordinates. Local geometric-engineering
work uses a protected notion of dynamical complex deformation that is checked
against Higgs chiral rings and magnetic quivers
\cite{Closset:2020afy,DeMarco:2026higgs}. Strict $L^2$ normalizability of a
representative three-form can additionally depend on the asymptotic metric of
the chosen noncompact Calabi--Yau, with subtleties already visible for the
conifold. We therefore separate:
\begin{enumerate}
 \item compact normalizability, which is sufficient for the theorem here;
 \item protected local Higgs dynamics, established for the node, $E_2$ and
 $E_3$ components used above; and
 \item strict local $L^2$ realization in a specified noncompact metric, which
 is not proved in this paper.
\end{enumerate}

No K\"ahler scaling ray has yet been constructed for $X_9$ or $X^{\circ}$
which determines exactly which shared rows survive after gravity decouples.
Accordingly, our principal theorem concerns compact M-theory at finite volume.
Claims about a particular rigid limit must use
\eqref{eq:projected-decoupling-matrix} and supply that additional geometric
analysis.
